\section{The concentration path through the actual Schur conductor map}
\label{sec:ns-schur-path}

This section composes the previously proved conic path with the integral
Schur endpoint, retaining the original source and target bases. Put
\[
 I=(0,2\sin^2(\pi/7)),\qquad
 B(\tau)=-42+196\tau-280\tau^2+160\tau^3-32\tau^4,
\]
and use exactly the sheet in \eqref{eq:ns-unit-circle-growth-map}.
The concentration parameter is $Q=\tau$, whereas the GCT parameter is
$q=e^{i\arcsin\sqrt{\tau/2}}$. These parameters remain distinct.
The endpoint factorization of Section~\ref{sec:plethysm-endpoint-transport}
pulls back to the three based maps
\begin{align}
 D_\tau&=\operatorname{diag}(1,1,1,\tau B(\tau)),\notag\\
 H_\tau&=\operatorname{diag}(1,1,1,h_\tau),\qquad
 V_\tau=\operatorname{diag}(1,1,1,-2\tau),\notag\\
 h_\tau&=-B(\tau)/2
 =21-98\tau+140\tau^2-80\tau^3+16\tau^4,\qquad
 D_\tau=V_\tau H_\tau.\label{eq:ns-schur-original-factorization}
\end{align}
Here $h_\tau$ denotes the retained conductor factor $[7]_q[3]_q$,
and is not the source profile parameter $h$. Indeed $r^2=-2\tau$
and $C=r^2[7]_q[3]_q=\tau B(\tau)$ prove every entry.
On $I$ both $C$ and $r^2$ are negative, so $h_\tau>0$.
The pointwise maps have their earlier endpoint, intermediate, and node modules
as source and target after this base change; the displayed matrices
do not identify those three original modules.

For the formal series statement below, the exact ring map must also be
specified. Set $R_m=\mathbb Z[m]$ and
\[
 h_m=(m+3)\bigl((m+4)^3-5(m+4)^2+6(m+4)-1\bigr),
 \qquad C_m=mh_m.
\]
Take three free $R_m$-modules with the respective original endpoint,
intermediate, and node basis labels. Define their maps to have matrices
$\operatorname{diag}(1,1,1,h_m)$ and
$\operatorname{diag}(1,1,1,m)$. Their composition has fourth entry
$C_m$. The homomorphism $R_m\to R=\mathbb Z[1/2][q,q^{-1}]$ sending
$m$ to $(q-q^{-1})^2$ takes $m+4$ to $(q+q^{-1})^2$ and hence
$h_m$ to $[7]_q[3]_q$. Mapping each free basis symbol to its stated
original basis vector is an isomorphism after this base change: the
inverse sends the original coordinate tuple back to the same symbols.
The three matrices show that these isomorphisms intertwine every map.
The same argument holds after the localizations used earlier.
This proves descent of this based endpoint diagram to $R_m$.
Its separate base change $R_m\to\mathbb R[[\tau]]$, $m\mapsto-2\tau$,
is the formal diagram used here. In particular, no homomorphism sending
the full Laurent variable $q$ to an element of $\mathbb R[[\tau]]$ is
asserted: the pointwise sheet has a square-root term at $\tau=0$.
The complete Schur base-change proof carries these diagram
isomorphisms through $F$ without changing any tableau labels.

Let $F=S_\lambda$, or $F=S_\lambda\circ S_\mu$, be the precise
integral Schur functor constructed above. Write $d=|\lambda|$ in
the first case and $d=|\lambda||\mu|$ in the second. Its tableau
weight multiplicities are the previously constructed $m_F(k)$.
Functoriality and the full diagonal formula give
\begin{equation}\label{eq:ns-schur-path-diagonal}
 \begin{aligned}
 F(D_\tau)&=\operatorname{diag}_{k,j}
       \bigl(\tau^k B(\tau)^k\bigr),\\
 F(H_\tau)&=\operatorname{diag}_{k,j}(h_\tau^k),\qquad
 F(V_\tau)=\operatorname{diag}_{k,j}((-2\tau)^k),\\
 &1\leq j\leq m_F(k).
 \end{aligned}
\end{equation}
For a single tableau, $k$ is its number of fourth symbols. In the
nested case, $k$ is the sum of those counts over the inner tableaux
appearing in the outer tableau, with multiplicity. Substitution of the fourth diagonal entry therefore
gives its $k$th power; the other three diagonal entries are one.
This also proves the displayed composition entry by entry.

Over $\mathbb R[[\tau]]$, $B$ is a unit because its constant term
is $-42$. An inverse is obtained recursively: if
$B=\sum_{i=0}^4b_i\tau^i$, set $a_0=-1/42$ and
$a_n=-b_0^{-1}\sum_{i=1}^{\min(4,n)}b_i a_{n-i}$.
Thus the exact image in the $(k,j)$ target coordinate is
$\tau^k\mathbb R[[\tau]]$, its kernel is zero, and its cokernel
is $\mathbb R[[\tau]]/(\tau^k)$. The latter is zero when $k=0$;
for $k>0$ its basis over $\mathbb R$ is $1,\tau,\ldots,\tau^{k-1}$,
as follows by separating the first $k$ coefficients of a series.
Consequently the cokernel length and determinant order both equal
$\sum_k k m_F(k)$. Specialization at $\tau=0$ has rank $m_F(0)$
and kernel and cokernel dimension $\sum_{k>0}m_F(k)$.
These are statements over the specified real series ring. The integer
factors $21^k$ and residual characteristics in the earlier integral
calculation are retained there; they have not been replaced by assertions
about valuations of prime numbers along a real path.

For each $\tau\in I$, identify source and target tableau labels by
their explicitly matching based-vector-space map and give the source
the Hermitian metric for which that basis is orthonormal. The form
defined by $F(D_\tau)$ then has coefficients $\tau^kB(\tau)^k$.
Since $B(\tau)<0$ on $I$, the even-weight and odd-weight coordinate
spaces are respectively positive and negative definite. They are
orthogonal and exhaust the space. A larger positive subspace would
intersect the negative coordinate space by the dimension formula,
and conversely for a larger negative subspace. Thus throughout $I$
the exact inertia is
\[
 \left(\sum_{k\ \text{even}}m_F(k),
       \sum_{k\ \text{odd}}m_F(k),0\right).
\]
The factor $F(H_\tau)$ is positive definite in the matched metric,
so the same indices hold for $F(V_\tau)$.

For the original shape $\nu=(7,5,3)$, the completely evaluated profile is
\[
 (m_\nu(0),\ldots,m_\nu(7))
 =(27,81,162,270,300,252,126,42).
\]
Its sum is $1260$ and its weighted sum is $4725$. It follows that
\begin{equation}\label{eq:ns-schur-source-path-data}
 \det S_\nu(D_\tau)=\tau^{4725}B(\tau)^{4725},\qquad
 \operatorname{inertia}=(615,645,0),\qquad
 \operatorname{rank}_{\tau=0}=27.
\end{equation}
The specialized kernel and cokernel have dimension $1233$; the series
cokernel has length $4725$. The endpoint entries extend to $\tau=0$,
although the conic connection does not extend across its excluded
branch locus. No connection value at that branch point is used here.

There is a second exact composition that tests what happens to growth.
Use the actual scalar functions of \eqref{eq:ns-negative-germ-transport}:
fix $0<h<1/100$, $A=1/2+h$, $e_0>0$, and bounded real $\varepsilon$,
and set $g(\tau)=\tau^{-A}(e_0+\tau^{2h}\varepsilon(\tau))$.
Define the explicit new linear map
\[
 L_\tau=g(\tau)D_\tau
\]
between the same two based endpoint spaces. This definition attaches
the stated scalar germ to their map; it does not identify an endpoint
vector with a fluid velocity or a quantum state. Each tableau monomial
has total degree $d$, including every inner and outer box in the nested
case, whence
\begin{equation}\label{eq:ns-schur-growth-exponents}
 F(L_\tau)_{k,j}
 =g(\tau)^d C(q)^k
 =\tau^{k-dA}B(\tau)^k
       (e_0+\tau^{2h}\varepsilon(\tau))^d.
\end{equation}
This proves homogeneity without a division by a factorial. For fixed
$k,d$, the binomial formula and boundedness give
$(e_0+\tau^{2h}\varepsilon)^d=e_0^d+O(\tau^{2h})$.
Also $B(\tau)^k=(-42)^k+O(\tau)$ for $k>0$, and the formula is
exactly one for $k=0$. Since $2h<1$, multiplication yields
\[
 F(L_\tau)_{k,j}
 =(-42)^k e_0^d\tau^{k-dA}
       +O(\tau^{k-dA+2h}).
\]
The leading constant is nonzero. Hence the absolute value tends to
zero when $k>dA$, tends to the absolute leading constant when $k=dA$,
and diverges when $k<dA$. This classification proves the outcome for
the defined map, rather than assuming that scalar multiplication and
a degree-$d$ functor have the same effect on an exponent.

In particular $d=15$ and $0\leq k\leq7$ for $\nu=(7,5,3)$, so
$k-15A\leq-1/2-15h<0$. Every one of its $1260$ diagonal entries
diverges in absolute value, with $615$ positive and $645$ negative
entries for sufficiently small $\tau$. Its determinant is exactly
\[
 \det S_\nu(L_\tau)
 =\tau^{4725-18900A}B(\tau)^{4725}
       (e_0+\tau^{2h}\varepsilon(\tau))^{18900},
\]
with leading exponent $-4725-18900h$. In contrast, the already proved
single scalar product $C(q)g(\tau)$ tends to zero with exponent
$1/2-h$. The displayed functorial computation explains this difference:
the degree-fifteen scalar contributes exponent $-15A$, and the largest
available fourth-coordinate weight contributes only $7$.

All multiplicities here are the unchanged nonnegative tableau counts.
The rational source comparison at $q=1$ is the explicit map already
proved in Section~\ref{sec:plethysm-endpoint-transport}; this path
calculation does not create a generic quantum-group intertwiner or an
unrestricted complexity lower bound. It supplies the exact Schur
composition, signs, and growth exponents on the specified original
endpoint instead.
